% called by main.tex
%
\chapter{Conclusiones}
\label{ch::conclusiones}

Este trabajo se planteó averiguar si el aprendizaje profundo puede anticipar el movimiento
de corto plazo del precio de los contratos binarios sobre Bitcoin de Polymarket y si esa
capacidad, en caso de existir, sobrevive a los costes y a la latencia de operar. La
respuesta a la primera pregunta es que sí, parcialmente y con matices; la respuesta a la
segunda es que no. Lo que da cuerpo a la memoria es el detalle de cómo se ha establecido
cada una de esas dos respuestas.

\section{La conclusión central, alcanzada por cuatro caminos}

Una precisión direccional cercana al 90\,\% no garantiza una política económica. Esta
afirmación, que vertebra el trabajo, no descansa sobre una única medición, y su solidez
proviene precisamente de que cuatro vías independientes convergen en ella.

La primera es la latencia. Medida como \emph{markout} desde el instante exacto de la señal,
la política produce $+6{,}79$ ticks de media; en el primer escalón que la malla de captura
permite resolver —dos segundos— el resultado ya es negativo. Conviene ser preciso sobre el
alcance de este argumento, y el Capítulo~\ref{ch::metodos} lo es: la latencia de ejecución
se midió contra el mercado real en unos $430$\,ms, incluida una orden efectivamente casada,
mientras que el escalón de dos segundos es la resolución del instrumento y no una latencia
observada. Entre ambos valores no existe ninguna observación, de modo que lo demostrado es
que la ventaja no sobrevive a dos segundos, no que muera a los cuatrocientos milisegundos.

La segunda vía es la estructura de comisiones, y es independiente de la anterior. La
auditoría del modelo de costes hasta las fuentes primarias reveló un error en la tarifa
empleada. Corregido y documentado con marca de tiempo verificable, el resultado es
inequívoco: con la tarifa oficial vigente, la comisión media de entrada agota por sí sola
el resultado bruto de la señal, sin que la latencia intervenga en el argumento. Ninguna
política del lado agresivo del libro es viable bajo ese régimen de costes.

La tercera es la ausencia de retardo explotable en la incorporación de información externa.
Si el contrato tardase varios segundos en reflejar los movimientos del subyacente, una
latencia inferior al segundo dejaría margen. La correlación cruzada entre el retorno del
futuro perpetuo de Bitcoin y la variación posterior del precio del contrato, medida sobre
841~mercados independientes y quince desfases, no supera $0{,}004$ en ningún caso. El
mercado digiere la información más rápido de lo que este montaje permite observar.

La cuarta es la única que llegó por validación prospectiva del lado \emph{taker}, y es la
que cierra el argumento con la potencia que a las demás les faltaba. El corte del protocolo general se ejecutó dos veces, cada una sobre la ventana
sellada de antemano: la primera hasta el 23 de agosto, sobre cincuenta días y 25.011
unidades de acción; la segunda hasta el 28, sobre cincuenta y cinco días y 27.208. La capa
histórica queda en $-0{,}004$ y en $+0{,}075$~ticks respectivamente, con un intervalo que
contiene el cero en los dos casos y la deja por tanto indistinguible del azar. Las capas
que cobran la comisión oficial caen entre $-0{,}69$ y $-3{,}65$~ticks según el corte, con
intervalos que lo excluyen holgadamente en ambos. Su valor no
está solo en el signo. Al calcular el intervalo de la cifra de referencia contra la que
se iba a comparar, los $+0{,}349$~ticks de la validación de junio, resultó que
\textbf{aquella nunca había sido significativa}: $[-0{,}194;\ +0{,}892]$ por filas y
$[-0{,}655;\ +1{,}337]$ por mercado. Entre junio y agosto, por tanto, no hay deterioro que explicar: hay una ausencia que
antes no podía afirmarse y ahora sí, con más de treinta veces más datos detrás. Lo que el corte aporta es lo que a la validación
original le faltaba: potencia suficiente para afirmar la ausencia en lugar de sugerirla.

\section{Qué hizo y qué no hizo el aprendizaje profundo}

La progresión de modelos profundos —arquitecturas recurrentes y convolucionales sobre el
libro, auto-atención, y un modelo fundacional de series temporales— aprende
representaciones útiles pero no produce una política estable. El modelo fundacional mejora
de forma monótona con la adaptación y aun así no alcanza a la deriva ingenua ni a un ARIMA,
y la curva de escalado de datos es plana.

Esa conclusión se sometió después a una prueba más exigente que la original, motivada por
una limitación detectada en el propio diseño: la curva publicada hacía crecer el número de
ventanas de entrenamiento sobre un conjunto fijo de trayectorias, de modo que medía densidad
de muestreo y no diversidad. Ampliando el corpus a las tres temporalidades capturadas y
llevando el eje hasta agotar los datos disponibles —casi el doble de ventanas de las
exploradas originalmente y casi el triple de trayectorias— ninguna de las dos vías, tabular
o profunda, se despega del azar en la predicción del signo.

De ese experimento surgió además un hallazgo que no se buscaba y que acota el problema mejor
que el resultado principal: la definición de la tarea excluye estructuralmente a la mayoría
de los mercados capturados. Entre contexto, horizonte y margen de seguridad, cada serie
requiere del orden de 376~fotogramas, y la longitud mediana de un mercado de cinco minutos
es de 160. Lo que limita el aprendizaje aquí no es solo el régimen de los datos: es la
geometría de la tarea, que descarta el 94\,\% de los mercados antes de que empiece el
entrenamiento. Una continuación que quiera aprovechar los mercados de horizonte corto no
necesita más datos, necesita redefinir el horizonte de predicción.

Quedaba un eje sin explorar, y cerró el argumento. Todos los modelos anteriores entrenaron
contra pérdidas estadísticas y decidieron después con un umbral; la última línea
(adenda~D, Sección~\ref{sec::adenda_d}) entrenó por fin la decisión directamente contra el
objetivo económico, con el contrafactual exacto de cada acción posible y la propia puerta
de riesgo del pre-registro escrita como función de pérdida. El resultado completa la
respuesta sobre el aprendizaje profundo por una vía independiente de la arquitectura y de
la escala: la frontera riesgo--rendimiento aprendida queda dominada en sus dieciséis puntos
por la de un parámetro escalar bien elegido. El eje de la arquitectura se amplió después con
tres modelos más —una LSTM, que cierra la única ablación que faltaba, y las dos familias
dominantes hoy en predicción de series temporales— y con cinco semillas por variante: las tres
resultan indistinguibles del control tabular, y la mejor de ellas iguala su rendimiento con
\textbf{trece veces menos parámetros}. Tras agotar los tres ejes —arquitectura, datos y
objetivo—, la conclusión ya no es que falte capacidad, sino que en este mercado y a esta
resolución no hay estructura explotable que una red pueda capturar y un escalar no; que el
modelo más pequeño de todos los probados empate con los demás es la forma más directa de
decirlo.

\section{La validación prospectiva}

El giro al lado pasivo del libro —motivado por que la misma estructura de comisiones que
penaliza a quien cruza el \emph{spread} retribuye a quien provee liquidez— condujo a una
política candidata que se sometió a validación prospectiva pre-registrada.

El protocolo se selló criptográficamente antes de observar los datos, con las huellas de los
cinco artefactos ejecutables, y se ejecutó una sola vez sobre una ventana de 35~días fijada
en el propio sello —que es de la tarde del primer día de esa ventana— y no observada hasta
el corte. El resultado es \textbf{NO-GO}: la política acumula $+3.408$\,\$
simulados, pero su peor día desciende por debajo del umbral de riesgo declarado de antemano,
y la regla se aplica sin interpretación.

Conviene precisar qué demuestra y qué no. Lo que demuestra es disciplina: un criterio
pre-registrado que solo aprueba no prueba nada, y este rechazó una estrategia
económicamente favorable porque violaba una condición escrita antes de mirar. Lo que no
demuestra es que la política careciera de valor, y el análisis posterior obliga a decirlo
sin rodeos: la puerta estaba prácticamente condenada de antemano. Bajo la dispersión
efectivamente observada, el peor día esperado en una ventana de ese tamaño era de
$-232$\,\$ frente a un umbral de $-150$, de modo que la probabilidad de fallarla superaba
el 60\,\% \emph{con independencia de la calidad de la política}. Dicho de otro modo: el
sub-criterio que firmó el veredicto era el único de los tres que carecía de potencia para
distinguir una política buena de una mala, mientras el que sí la tenía —equity neta
positiva— se cumplía con holgura. El NO-GO es un resultado sobre el procedimiento, no una
medición del valor económico de la estrategia, y presentarlo como lo segundo sería
justamente el tipo de confusión que este trabajo intenta evitar.

Una circunstancia externa reforzó involuntariamente la validez del ejercicio: la ventana no
se cerró por decisión del autor sino por una desconexión física del almacenamiento del
recolector, tres días antes de que se calculara la primera métrica. El punto de corte, por
tanto, no pudo verse influido por los datos que iba a revelar.

\section{Aportaciones}

Las aportaciones de este trabajo son metodológicas antes que predictivas.

\begin{enumerate}[noitemsep]
  \item Un sistema de captura multifuente con control de calidad por sesión, operando de
  forma continua, y el corpus alineado que produce.
  \item Un protocolo de evaluación que combina validación estrictamente temporal, un modelo
  explícito de costes con la tarifa oficial y una latencia medida contra el mercado real,
  distinguiendo qué parte de ella es medición y cuál es resolución del instrumento.
  \item Evidencia cuantitativa, por cuatro caminos independientes, de la ruptura entre
  predicción y ejecución, incluida una validación prospectiva del lado agresivo con
  potencia suficiente para afirmar la ausencia de señal, y no solo para sugerirla.
  \item Una validación prospectiva pre-registrada y sellada, con resultado negativo
  reportado íntegro.
  \item Una trazabilidad que conserva los resultados negativos y las correcciones de error
  con la misma dignidad que los positivos, incluida una fe de erratas del modelo de costes
  documentada con marca de tiempo verificable.
\end{enumerate}

No se afirma rentabilidad real ni se presenta un sistema listo para desplegar.

\section{Lecciones metodológicas}

Un trabajo cuyo objeto es la honestidad estadística debe reportar también dónde falló su
propia disciplina. Cinco patrones se repitieron lo suficiente como para merecer mención.

El primero es la fijación de umbrales de decisión sin análisis de potencia. Ocurrió dos
veces: el umbral de peor día del pre-registro se declaró en cantidad absoluta, sin advertir
que codifica implícitamente una hipótesis sobre la dispersión; y el criterio de separación
entre arquitecturas se fijó por debajo del error estándar del propio conjunto de evaluación,
de modo que ninguna de sus dos ramas podía establecerse. En ambos casos el umbral se eligió
por analogía con una magnitud observada antes, en lugar de derivarlo de la dispersión del
estadístico que iba a contrastar. La lección es que el pre-registro es condición necesaria
para la honestidad de un experimento, pero no suficiente para su potencia.

El segundo es la confusión entre un oráculo y una política. En dos ocasiones un análisis
sugirió un margen de mejora considerable que, examinado, resultaba proceder de conocer por
anticipado qué operaciones perderían dinero y evitarlas. Un techo calculado con información
del futuro no acota lo que un modelo puede lograr; solo acota lo que lograría alguien que ya
conoce el resultado.

El tercero es que un resultado nulo puede estar midiendo el instrumento y no el fenómeno.
La primera frontera riesgo--rendimiento de la adenda~D devolvió dieciséis puntos idénticos
hasta el último decimal; la reacción correcta no fue aceptar el «no aporta» —que era la
conclusión cómoda y esperada— sino instrumentar el optimizador, que reveló \emph{logits}
saturados con gradiente del orden de $10^{-6}$: la política no elegía la constante, estaba
congelada en ella. Solo tras corregirlo y comprobar que la política se movía, el mismo
veredicto negativo pasó a ser evidencia. Un resultado demasiado exacto es una alarma, no
una confirmación.

El cuarto es que la unidad de análisis correcta depende de la pregunta. La regla de un
\emph{token} por mercado es obligada al contar muestras independientes en un contraste
estadístico, y esa misma regla, trasladada a una evaluación económica, fabricó un falso
cumplimiento de la puerta de riesgo: en la operativa se cotizan los dos \emph{tokens} y el
umbral se calibró sobre el libro completo, de modo que evaluar con medio libro reducía el
peor día por construcción. La regla no era incorrecta; estaba aplicada a la pregunta
equivocada.

El quinto afecta al cruce entre fuentes de datos. Los fallos de alineación temporal no
producen errores: producen tablas vacías, o casi vacías, que superficialmente parecen
resultados. En este trabajo un desajuste entre dos relojes con precisiones distintas
devolvió una tabla completa de correlaciones que no significaba nada, y su única señal de
alarma fue que los valores eran demasiado exactos. Todo cruce entre fuentes debe asertarse,
nunca suponerse.

\section{Trabajo futuro}

La continuación inmediata es la validación de las conclusiones de política sobre la ventana
de captura posterior al corte, que permanece sin observar y crece cada día.

Una segunda línea, más ambiciosa, ataca la limitación estructural del simulador maker —su
incapacidad para medir la respuesta competitiva del mercado (Capítulo~\ref{ch::etica})—. La
reproducción histórica trata el libro como un decorado inerte; capturar cómo reaccionarían
los creadores de mercado incumbentes a las órdenes propias exige un simulador
\emph{reactivo}. La línea de simuladores generativos de libros de órdenes de los últimos
años ofrece el instrumental: modelos de «agente mundo» que aprenden de datos históricos una
dinámica del libro capaz de reproducir cierto impacto de
mercado~\cite{coletta2022worldagent}, modelos generativos autorregresivos del flujo de
mensajes a nivel de \emph{token}~\cite{nagy2023generative} y, más recientemente, modelos de
difusión para la simulación y predicción del libro~\cite{backhouse2025diffusion}; una
revisión reciente sistematiza el campo~\cite{jain2024lobreview}. Integrar uno de estos
simuladores con la política congelada de este trabajo permitiría estimar la selección
adversa y el reparto de \emph{rebates} bajo competencia, y no bajo el supuesto optimista de
la reproducción histórica. Tiene además una motivación cuantificada: el supuesto de
ejecución sostiene cerca de dos tercios de las órdenes casadas del simulador y hoy no puede
contrastarse contra nada.

Finalmente, la limitación de resolución temporal del instrumento sugiere una tercera línea:
capturar el libro a una frecuencia sensiblemente mayor permitiría observar el intervalo
entre cero y dos segundos, donde vive la latencia de ejecución realmente medida y que este
trabajo no ha podido explorar. Esa línea exigiría además migrar el camino caliente de la
captura y de la ejecución de Python a un lenguaje compilado como C++ o Rust. La ventaja no
es genérica sino específica de este cuello: sin recolector de basura ni bloqueo global del
intérprete, las colas de mensajes se procesan con latencias del orden de microsegundos y,
sobre todo, \emph{predecibles} —en el camino medido aquí, la varianza de la latencia
importa tanto como su media—. No haría falta partir de cero: existen conectores y motores
de libro de órdenes maduros en ambos lenguajes (el ecosistema de \emph{trading}
cuantitativo de código abierto en Rust es hoy notablemente activo), de modo que el trabajo
propio se concentraría en el cliente de la API de Polymarket y en la lógica de decisión,
manteniendo en Python el análisis y el entrenamiento, donde la latencia no manda.

La regla que guió el proyecto permanece: la complejidad y cualquier afirmación económica
deben ganarse su sitio con evidencia verdaderamente nueva.
